% ── Section 8: Conclusion ─────────────────────────────────────────────────────
\section{Conclusion}
\label{sec:conclusion}

This work presented \textbf{BanglaMix}, the first large-scale dataset and
benchmark for dialectal Bengali--English code-switching ASR.
The dataset comprises 51,857 clips totalling 164.74~hours across 15 regional
dialects of Bangladesh, collected from YouTube and Facebook using
dialect-name search queries targeting content genres where natural CS is
most prevalent.
An automated pipeline based on WebRTC VAD segmentation, Whisper
silver-standard transcription, and a single-model LLM judge
(Qwen2.5) for borderline clip validation enables large-scale annotation
without any human transcribers. For evaluation, a separate three-model LLM
ensemble (Mistral-7B, LLaMA-3.1-8B, Gemma) rates hypotheses beyond surface WER.

This work benchmarked three ASR architectures---Whisper, Meta MMS, and
wav2vec2-XLS-R---under zero-shot and fine-tuned conditions.
Key findings are:

\begin{enumerate}[leftmargin=2em, itemsep=2pt]
  \item Zero-shot multilingual models substantially underperform on dialectal CS
        speech ($\geq 0.91$~WER), even state-of-the-art Whisper large-v3.
  \item Dialect-only fine-tuning (Ablation A) and code-switch-only fine-tuning
        (Ablation B) trade off Bengali vs.\ English recognition: BN\_ONLY barely
        reproduces English (EN survival 0.18), while CS\_ONLY best preserves it
        (0.47) at the cost of overall accuracy. Training on \emph{all} data gives
        the best balance.
  \item The proposed Whisper FT ALL model achieves \textbf{0.80~WER /
        0.49~CER} overall (0.82~WER on CS clips), a 25\% relative improvement
        over its zero-shot counterpart and a statistically significant gain over
        every baseline ($p<0.001$).
  \item Absolute error remains high: dialectal Bengali has no standardised
        orthography and labels are silver-standard, so BanglaMix is best viewed
        as a challenging open benchmark rather than a solved task.
  \item Per-dialect analysis reveals that the most phonologically distinct
        dialect, Sylhet, is the hardest across all models, suggesting future
        work should prioritise targeted data collection for such varieties.
\end{enumerate}

\paragraph{Limitations.}
BanglaMix relies on silver-standard transcripts from Whisper, which may
contain systematic errors---particularly for words spoken in strong dialect
with no standard Bengali equivalent.
The code-switched proportion ($\sim$7\%) reflects the natural prevalence of
CS in our collected content; more aggressive query targeting or domain
selection could increase this.
The dataset is limited to audio from social media, which may not generalise
to formal domains (news, education, healthcare).

\paragraph{Future work.}
Future directions include: (1)~expanding to dialects with no current
coverage (Jessore, Faridpur, Cox's Bazar); (2)~domain extension to medical
and academic code-switching; (3)~developing dialect identification as an
auxiliary task to improve ASR; (4)~human post-correction of the highest-impact
clips using active learning to maximise annotation efficiency.

\paragraph{Data and code availability.}
The BanglaMix dataset, all model checkpoints, and the complete data collection
and evaluation pipeline are released at \texttt{[repository URL]} under a
CC-BY 4.0 licence (audio from publicly available sources; transcripts under
CC0).

\section*{Acknowledgements}

[Acknowledgements will be added here.]
This work was supported by [funding agency / grant number].
